\documentclass{report}
\usepackage{amsmath,amssymb}
\setlength{\parindent}{0mm}
\setlength{\parskip}{1em}
\begin{document}
\begin{center}
\rule{6in}{1pt} \
{\large Eric C. Cyr \\
{\bf Spatially Varying Stochastic Expansions for Embedded Uncertainty Quantification}}

Sandia National Laboratories \\ P O Box 5800 \\ MS 1320 \\ Albuquerque \\ NM 87185-1320
\\
{\tt eccyr@sandia.gov}\end{center}

Uncertainty quantification (UQ) is increasingly recognized as necessary to assess
the effect of variability of input parameters on a physical model. For the solution
of stochastic PDEs, existing discretizations are based on a tensor product between
the deterministic basis and the stochastic basis. The implicit assumption
made by the tensor product is that the required order of the stochastic approximation
is fixed throughout the domain. However, solutions to many
PDEs of interest exhibit spatially localized features that may result in
uncertainty being severely over or under resolved by existing discretizations.
The result of under resolution is a loss in accuracy, and for over
resolution excess computational cost is incurred.

To address this issue, we will consider using a discretization where the order
of the stochastic expansion is allowed to vary across the physical
domain. This is achieved by a Galerkin projection onto a spatially varying
basis. The challenge of this method is in the construction of such a basis which,
due to the fine level of granularity, compels us to use
embedded UQ techniques that require a direct modification of the PDE code.
This is in contrast to black-box techniques, like Monte-Carlo or stochastic collocation,
which do not appear to be applicable to the spatially varying discretization.

In this talk, we propose an adaptive algorithm using numerical error indicators
to construct the spatially varying discretization.
Furthermore, we will discuss the changes to the stochastic weak form that make
the spatially varying discretization possible. Finally, we will demonstrate the
performance of the adaptive algorithm with several numerical
examples and provide comparisons to other embedded UQ techniques.


\end{document}
